\documentclass[11pt]{article}
\usepackage[margin=1in]{geometry}
\usepackage{amsmath,amssymb,amsthm,bm}
\usepackage{mathtools}
\usepackage{hyperref}
\usepackage{enumitem}
\usepackage{booktabs}
\usepackage{graphicx}
\usepackage{xcolor}

% ---- macros -------------------------------------------------------------
\newcommand{\R}{\mathbb{R}}
\newcommand{\T}{^{\!\top}}
\newcommand{\inv}{^{-1}}
\newcommand{\norm}[1]{\left\|#1\right\|}
\newcommand{\cN}{\mathcal{N}}
\newcommand{\cL}{\mathcal{L}}
\newcommand{\cP}{\mathcal{P}}
\newcommand{\cGP}{\mathcal{GP}}
\newcommand{\Omegabd}{\partial\Omega}
\newcommand{\Th}{\Theta}
\newcommand{\Thtr}{\Th_{\mathrm{train}}}
\newcommand{\eps}{\varepsilon}
\newcommand{\Ed}{E_d}
\newcommand{\Lw}{\mathcal{L}_{\bar w}}
\newcommand{\bigO}{\mathcal{O}}

\title{A near-linear sparse-Cholesky solver for the Darcy inverse problem}
\author{Working note}
\date{\today}

\emergencystretch=3em

\begin{document}
\maketitle

\begin{abstract}
We give a Gaussian-process / PDE-collocation solver for the Darcy inverse
problem --- recover both the diffusion coefficient $a=e^{w}$ and the state $u$
from a forcing, boundary data, and noisy interior measurements of $u$ --- whose
total cost grows essentially \emph{linearly} in the number of collocation points
$N$, matching the efficiency of a single forward elliptic solve up to a small
constant factor (we solve for two unknown fields rather than one). The solver
uses two ingredients, both built only from the near-linear primitives supplied
by an incomplete sparse-Cholesky (Kullback--Leibler, ``KL'') factorization of
the \emph{covariance} matrices of two \emph{independent} priors $K_u$ (for $u$)
and $K_w$ (for $w$); the two kernels are never summed and the precision Hessian
is never assembled. (i) An \emph{alternating refinement} stage in which each
half-step is a linear GP regression. The hard half --- the $u$-update with noisy
interior data co-located on the PDE collocation points --- is solved by a direct
\emph{bordered / Schur} solve whose only non-forward ingredient, the data-block
posterior covariance, is read off in closed form (and in near-linear time) from
the trailing block of the joint screened precision factor. (ii) A
\emph{warm-up} stage that finds the basin of the (non-convex) MAP problem by
running a few joint Gauss--Newton steps on a \emph{fixed coarse subset} of the
points and then kriging the recovered coefficient field onto all $N$ points; its
cost is $\bigO(1)$ in the fine $N$. We report controlled scaling experiments
from $N=800$ to $N=12800$ showing a flat per-stage profile and an
overall empirical exponent $\approx 1$, with coefficient-recovery error that is
robust in $N$.
\end{abstract}

\tableofcontents

% =========================================================================
\section{Problem and scope}\label{sec:problem}
% =========================================================================

We consider the Darcy (steady diffusion) inverse problem on $\Omega=[0,1]^2$:
recover the diffusion coefficient $a(x)=e^{w(x)}>0$ and the state $u(x)$ from
\begin{equation}\label{eq:darcy}
\begin{aligned}
-\nabla\!\cdot\!\bigl(e^{w(x)}\nabla u(x)\bigr) &= f(x), && x\in\Omega,\\
u(x) &= g(x), && x\in\Omegabd,\\
y_j &= u(x_j^{d}) + \eta_j, && \eta_j\sim\cN(0,\sigma^2),\ \ j=1,\dots,N_d.
\end{aligned}
\end{equation}
The forcing $f$, boundary data $g$, observation sites $x_j^d$, observations
$y_j$, and noise level $\sigma\ge 0$ are given. We parameterize $a=e^w$ so that
positivity is automatic and $w=\log a$ is unconstrained. Expanding the
divergence,
\begin{equation}\label{eq:pde-expanded}
-e^{w}\bigl(\Delta u + \partial_1 w\,\partial_1 u + \partial_2 w\,\partial_2 u\bigr)=f
\quad\Longleftrightarrow\quad
\Delta u = -\,\partial_1 w\,\partial_1 u - \partial_2 w\,\partial_2 u - f\,e^{-w}.
\end{equation}

\paragraph{Test problem.} Throughout we use $f\equiv 1$, $g\equiv 0$, and the
two-peak truth
\begin{equation}
a_\star(x) = e^{\sin 2\pi x_1+\sin 2\pi x_2}
           + e^{-(\sin 2\pi x_1+\sin 2\pi x_2)}
         = 2\cosh\!\bigl(\sin 2\pi x_1+\sin 2\pi x_2\bigr),
\end{equation}
i.e.\ $w_\star=\log a_\star$. The reference state $u_\star$ is produced by a
fine finite-difference forward solve; data $y_j$ are $u_\star$ sampled at $N_d$
interior points plus Gaussian noise of level $\sigma=10^{-3}$. Recovering the
two peaks of $a_\star$ is the qualitative success criterion.

\paragraph{Goal.} Solve \eqref{eq:darcy} with a cost that is \emph{near-linear}
in $N$ --- the same complexity class as solving a single (forward) elliptic
equation with the sparse-Cholesky machinery --- allowing only a constant-factor
overhead for the second unknown field. We restrict ourselves to the primitives
of the KL sparse-Cholesky factorization and explicitly \emph{exclude} external
fast methods (multigrid, nested dissection, FMM, $H$-matrices). The central
difficulty, as in the forward problem, is to solve the arising linear systems
with good (screening) preconditioners.

% =========================================================================
\section{GP collocation and the sparse-Cholesky primitives}\label{sec:gp}
% =========================================================================

\subsection{Latent fields, functionals, and the MAP problem}

We place \emph{independent} mean-zero Gaussian-process priors
$u\sim\cGP(0,K_u)$ and $w\sim\cGP(0,K_w)$, both with Mat\'ern-$7/2$ kernels but
with different length scales: a long $\ell_u=1.0$ (the state solves an elliptic
PDE and is smooth, so its derivatives $\nabla u,\Delta u$ --- which pin $a$
through \eqref{eq:pde-expanded} --- must be clean) and a short $\ell_w=0.2$ (the
coefficient has sharp peaks and the $w$-factor must stay local/sparse). Using a
single shared, short length scale flattens the recovered $a$; this length-scale
split is essential for two-mode recovery.

We discretize each field by its \emph{jet} at $N$ interior collocation points
$\{x_i\}_{i=1}^N\subset\Omega$ and the boundary points
$\{x_i^\partial\}_{i=1}^{N_\partial}$. Collect the unknowns
\begin{equation}
z = \bigl(\,\underbrace{w,\,\partial_1 w,\,\partial_2 w}_{w\text{-jet}},\;
            \underbrace{u,\,\partial_1 u,\,\partial_2 u}_{u\text{-jet}}\,\bigr)\in\R^{6N},
\end{equation}
evaluated at the $x_i$ (the Laplacian $\Delta u$ is \emph{not} a free unknown:
it is eliminated pointwise through the PDE, see \eqref{eq:v3} below). Each field
is represented in the GP by a family of linear functionals --- point value
$\delta_x$, first derivatives $\partial_1,\partial_2$, and Laplacian $\Delta$
--- applied at the collocation/boundary points. For a stacked set of $M$ such
functionals $L=(L_1,\dots,L_M)$ the prior Gram matrix is
$\Theta = \bigl(\mathrm{Cov}[L_a\phi,\,L_b\phi]\bigr)_{a,b}=\bigl(L_a L_b
K\bigr)_{a,b}$, a single-kernel covariance.

The MAP estimate minimizes the sum of the two prior RKHS norms and the data
misfit, subject to the boundary condition and the (nonlinear) PDE link. Writing
$v$ for the vector of $u$-functional values implied by $z$ and $w$ for the
vector of $w$-functional values, and using that for a GP the RKHS norm of a
function interpolating values $\xi$ on functionals $L$ equals $\xi\T\Theta\inv
\xi$, the objective is
\begin{equation}\label{eq:loss}
J(z) \;=\; v\T \Th_u\inv v \;+\; w\T \Th_w\inv w
        \;+\; \frac{1}{\sigma^2}\sum_{j=1}^{N_d}\bigl(u(x_j^d)-y_j\bigr)^2,
\end{equation}
where $\Th_u$ is the covariance of the $u$-functionals $\{\delta_\partial,
\delta_\Omega,\partial_1,\partial_2,\Delta\}$ and $\Th_w$ that of the
$w$-functionals $\{\delta_\Omega,\partial_1,\partial_2\}$. The boundary rows of
$v$ are fixed to $g$, and the Laplacian rows are eliminated through the PDE
\eqref{eq:pde-expanded}:
\begin{equation}\label{eq:v3}
(\Delta u)_i = -(\partial_1 u)_i (\partial_1 w)_i
              -(\partial_2 u)_i (\partial_2 w)_i - f_i\, e^{-w_i}.
\end{equation}
\textbf{The two priors enter \eqref{eq:loss} as two separate quadratic forms;
$K_u$ and $K_w$ are never added.} This is what keeps each Gram matrix a
single-kernel covariance that the screening theory can factor.

\subsection{The sparse-Cholesky primitive}\label{sec:kl}

The Schäfer--Sullivan--Owhadi screening theory shows that, in a fine-to-coarse
\emph{maximin} ordering and with a sparsity radius $\rho$ (keep, for each point,
the $\bigO(\rho^d)$ nearest already-ordered points), the Cholesky factor $U$ of
the \emph{inverse} covariance is exponentially accurate when truncated to that
sparsity pattern. The KL factorization returns a sparse upper-triangular $U$ and
a permutation $P$ with
\begin{equation}\label{eq:kl}
\Theta\inv \;\approx\; P\T (U U\T) P ,
\qquad \mathrm{nnz}(U)=\bigO(N\rho^d),
\end{equation}
both built in near-linear time. From \eqref{eq:kl} we get the only two linear
primitives we will use, each near-linear:
\begin{equation}\label{eq:prims}
\begin{aligned}
\Theta\inv x &\;=\; P\T U U\T P x &&(\text{two sparse mat-vecs}),\\
\Theta x &\;=\; P\T U^{-\!\top} U\inv P x &&(\text{two sparse triangular solves}).
\end{aligned}
\end{equation}
We build one such ``big'' factor for $\Th_u$ and one for $\Th_w$ once; the
``small'' factors used as preconditioners below are rebuilt per linear solve but
are themselves near-linear. We use $\rho=3$ throughout (the elliptic regime;
the forward solver of the companion note is $N$-independent at this $\rho$).

% =========================================================================
\section{The solver}\label{sec:solver}
% =========================================================================

The objective \eqref{eq:loss} is non-convex (through \eqref{eq:v3}). The solver
has two stages: a \emph{warm-up} that lands in the recovery basin
(Sec.~\ref{sec:warmup}) and an \emph{alternating refinement} that converges to
the MAP estimate (Sec.~\ref{sec:alt}). We describe the alternating stage first
because the warm-up is engineered around it.

\subsection{Alternating refinement}\label{sec:alt}

With one field frozen, \eqref{eq:loss} is a \emph{linear} GP regression in the
other. Each outer step performs a $u$-update then a $w$-update; both re-solve
their field from scratch, so only a coefficient field $w$ in the right basin is
needed to seed the stage.

\paragraph{$w$-update (easy).} With $u$ fixed, \eqref{eq:pde-expanded} is linear
in $w$ and the regression has a single functional family on $\Th_w$; it is small
and well-conditioned, converging in $\bigO(10)$ preconditioned-CG (pCG)
iterations with the $\Th_w$ small factor as preconditioner. It is never the
bottleneck.

\paragraph{$u$-update (the hard half).} Fix $w$ and write $\bar w_1=\partial_1
w$, $\bar w_2=\partial_2 w$. The frozen PDE is the linear operator
\begin{equation}
\Lw u := \Delta u + \bar w_1\,\partial_1 u + \bar w_2\,\partial_2 u,
\qquad \Lw u = -f e^{-\bar w} =: b^\Omega .
\end{equation}
The $u$-update is the posterior mean of $u\sim\cGP(0,K_u)$ under
\begin{equation}\label{eq:uconstraints}
\underbrace{u(x_i^\partial)=g_i}_{N_\partial},\qquad
\underbrace{\Lw u(x_i^\Omega)=b_i^\Omega}_{N_\Omega=N},\qquad
\underbrace{u(x_j^d)\approx y_j}_{N_d,\ \text{noise }\sigma}.
\end{equation}
In covariance (``$\alpha$-space'') form, with $\Thtr$ the Gram of these
$N_\partial+N+N_d$ functionals and $R=\sigma^2$ on the data rows,
\begin{equation}\label{eq:alpha}
(\Thtr + R)\,\alpha = \begin{pmatrix}g\\ b^\Omega\\ y\end{pmatrix},
\qquad
u(x_*) = k(\delta_{x_*},\text{train})\,\alpha .
\end{equation}

\textbf{Why a plain preconditioner fails, and the fix.} Without data
($N_d=0$) the system \eqref{eq:alpha} is exactly a forward elliptic solve: the
$\Th_u$ small factor preconditions it to an $N$-independent iteration count
($\sim 13$ at $\rho=3$ in our runs), reproducing the forward solver. The trouble
is that the data sites are \emph{co-located} with PDE points
($x_j^d=x_j^\Omega$) and $\sigma$ is tiny, so the data $\delta$-rows are nearly
redundant with the PDE rows: their conditional variance given $\{$bdy,
PDE$\}$ is $\approx\sigma^2 I$, a tight cluster of near-null Gram modes that a
local screened factor of the \emph{full} augmented operator cannot resolve at
$\rho\le 5$. (Raising $\rho$ patches one $N$ but the stiff-mode count grows
$\propto N_d$.)

The cure is to \emph{peel off} the small data block instead of preconditioning
through it. Partition the functionals into the forward block
$F=\{\text{bdy},\text{PDE}\}$ (which screens perfectly) and the data block
$d=\{N_d$ noisy $\delta\}$. With
$\Thtr=\begin{psmallmatrix}\Th_{FF}&\Th_{Fd}\\ \Th_{dF}&\Th_{dd}\end{psmallmatrix}$
the block solve of \eqref{eq:alpha} is exact:
\begin{equation}\label{eq:bordered}
\begin{aligned}
S &= S_0 + \sigma^2 I,\qquad S_0 = \Th_{dd} - \Th_{dF}\Th_{FF}\inv\Th_{Fd},\\
\alpha_d &= S\inv\bigl(y - \Th_{dF}\Th_{FF}\inv\, \text{rhs}_F\bigr),\qquad
\alpha_F = \Th_{FF}\inv\bigl(\text{rhs}_F - \Th_{Fd}\,\alpha_d\bigr).
\end{aligned}
\end{equation}
Here $\Th_{FF}\inv$ is applied by the forward pCG ($N$-independent iterations,
the $\Th_u$ small factor as preconditioner), $\Th_{Fd}\cdot$ and
$\Th_{dF}\cdot$ are single mat-vecs of the big-factor operator, and $S$ is a
small $N_d\times N_d$ dense system. The only delicate object is the data-block
posterior covariance $S_0$; forming it by $N_d$ forward solves would cost
$\bigO(N\cdot N_d)$. We avoid that entirely.

\paragraph{Near-linear data-block covariance via the trailing precision block.}
The block-inverse identity gives $S_0\inv=(\Thtr\inv)_{dd}$ exactly. The KL
factor of the \emph{joint} screened system over $[\text{bdy},\text{PDE},
\text{data}]$ already supplies $\Thtr\inv\approx P\T U U\T P$. Hence, with
$a_j$ the position of data index $d_j$ under the permutation $P$ (i.e.\
$a=P\inv[d]$),
\begin{equation}\label{eq:trailing}
(\Thtr\inv)_{d_i d_j} = (U U\T)_{a_i a_j} = U_{a_i,:}\,U_{a_j,:}\T,
\qquad\Longrightarrow\qquad
(\Thtr\inv)_{dd} = U_d\,U_d\T,
\end{equation}
where $U_d$ stacks the $N_d$ permuted data rows of $U$. So $S_0\inv$ is just the
dense Gram of $N_d$ sparse rows of $U$ --- formed in $\bigO(\mathrm{nnz}(U_d))$
time, \emph{no} forward solves. This is the key that turns the exact data-block
handling from $\bigO(N^2)$ into near-linear (up to the $\bigO(N_d^2)$ dense $S$,
discussed in Sec.~\ref{sec:complexity}). Numerically the screened $S_0$ carries
one spurious large-variance mode; it is benign, merely down-weighting one data
direction back to the accurate forward solve.

\subsection{Warm-up by a coarse solve and kriging}\label{sec:warmup}

The alternating stage converges only from inside the recovery basin: started
from $w=0$ it sits at the trivial fixed point (the $u$-update at $w=0$ forces
$\Delta u=-f$, killing the $w$-update's signal). A few \emph{joint}
Gauss--Newton (GN) steps on \eqref{eq:loss} escape it. The GN Hessian is the
precision/normal form
\begin{equation}\label{eq:H}
H = 2\,J_v\T \Th_u\inv J_v + 2\,J_w\T \Th_w\inv J_w + D,
\end{equation}
with $J_v,J_w$ the (sparse) Jacobians of the functional values w.r.t.\ $z$ and
$D$ the data diagonal. \emph{Assembling and factoring $H$ is the original
bottleneck.} The Jacobian sandwiching makes $H$ a squared (biharmonic-like)
operator: its exact Cholesky fills in ($\bigO(N^{2.8})$ per step), and --- unlike
the screened covariance precision of Sec.~\ref{sec:kl} --- its ill-conditioning
($\kappa(H)\sim10^{9}$) is \emph{not} of a form any incomplete factorization can
capture (no convergence at any $N$ or fill; we verify this in
Sec.~\ref{sec:verif} and distil the rule as Insight~6). So instead of factoring
$H$ we sidestep it.

The key observation is that the warm-up's \emph{only} product that the
alternating stage consumes is a coefficient field $w$ in the right basin
(everything else is recomputed in \eqref{eq:bordered}). So:
\begin{enumerate}[leftmargin=2em]
\item Run the joint GN warm-up on a \emph{fixed coarse subset} of $N_c$
  interior points (with $N_{d,c}$ coarse data points), at cost \emph{independent
  of the fine $N$}. We use $N_c=400$, $N_{d,c}=150$, $6$ GN steps.
\item Krige the coarse coefficient jet to all $N$ fine points. With $\alpha_w^c=
  (\Th_w^c)\inv w^c$ the coarse dual weights (one application of the coarse $w$
  factor), the fine field and its gradients are the GP posterior means
  \begin{equation}
  (w,\partial_1 w,\partial_2 w)(x_i)
   = k\bigl(\{\delta,\partial_1,\partial_2\}_{x_i},\ \text{train}_w^c\bigr)\,\alpha_w^c,
   \qquad i=1,\dots,N,
  \end{equation}
  one cross-covariance evaluation against the fixed coarse training set,
  $\bigO(N)$.
\end{enumerate}
This is \emph{not} multigrid: there is a single coarse solve and one kriging
evaluation, with no hierarchy, no cycles, and no restriction/prolongation
operators. The resulting fine $w$-field seeds the alternating stage. (We also
tested a ``data-first'' warm-up --- regress $u$ from data alone, then one
$w$-update --- but its data-only Laplacian is too noisy to seed a good $w$ and
it stalls; the coarse-solve warm-up is the method of choice.)

% =========================================================================
\section{Complexity}\label{sec:complexity}
% =========================================================================

Per-stage costs, with $\rho$ fixed and $N_d=\bigO(N)$:
\begin{itemize}[leftmargin=2em]
\item \textbf{Big factors} ($\Th_u,\Th_w$, built once): $\bigO(N\log N)$ build,
  $\bigO(N)$ apply.
\item \textbf{Warm-up:} coarse GN is $\bigO(1)$ in $N$; kriging is $\bigO(N)$.
\item \textbf{Alternating, per outer step:} small-factor builds $\bigO(N\log
  N)$; forward pCG $\bigO(N)\times(N\text{-independent iters})$; trailing-block
  $S_0\inv$ extraction \eqref{eq:trailing} $\bigO(\mathrm{nnz}(U_d))=\bigO(N)$;
  the dense $N_d\times N_d$ data solve $\bigO(N_d^2)$ to apply (a Cholesky
  factorization $\bigO(N_d^3)$).
\end{itemize}
Every ingredient is near-linear except the dense data block. With
$N_d\propto N$ that block is asymptotically $\bigO(N^2)$--$\bigO(N^3)$, but with
a tiny constant: over the tested range it is sub-dominant and the measured
profile is flat (Sec.~\ref{sec:results}). The known route to remove it ---
screen $S$ itself, since the data sites are spatially distributed so $S_0$
decays in space --- is left as future work; it does not affect the present
near-linear measurements.

% =========================================================================
\section{Numerical results}\label{sec:results}
% =========================================================================

All runs use Mat\'ern-$7/2$, $\ell_u=1.0$, $\ell_w=0.2$, $\rho=3$,
$\sigma=10^{-3}$, $6$ warm-up GN steps on a fixed $N_c=400$/$N_{d,c}=150$ coarse
subset, and $3$ alternating outer steps; single CPU thread (JAX CPU backend).
The coefficient error is $\mathrm{rel}_a=\norm{a-a_\star}_{L^2}/\norm{a_\star}_{L^2}$.

\begin{table}[h]
\centering
\caption{End-to-end scaling with the coarse-warm-up default and the near-linear
trailing-block $u$-step. Stage-1 (warm-up) is flat in $N$; total cost grows
$\approx$ linearly; coefficient recovery is $N$-robust.}
\label{tab:scaling}
\begin{tabular}{rrrrrr}
\toprule
$N$ & $N_d$ & Stage-1 (s) & Stage-2 (s) & Total (s) & $\mathrm{rel}_a$ \\
\midrule
$800$   & $300$  & $6.6$  & $9.4$   & $22.7$  & $7.89\%$ \\
$1600$  & $600$  & $6.9$  & $22.1$  & $41.0$  & $7.81\%$ \\
$3200$  & $1200$ & $7.5$  & $69.6$  & $99.5$  & $7.73\%$ \\
$6400$  & $2400$ & $11.3$ & $159.8$ & $214.9$ & $7.62\%$ \\
$12800$ & $4800$ & $16.0$ & $349.8$ & $448.9$ & $7.40\%$ \\
\bottomrule
\end{tabular}
\end{table}

\noindent\textbf{Reading the table.} (i) The warm-up (Stage 1) is flat ---
the joint GN runs on the same fixed coarse subset regardless of $N$, and the
only growing piece is the $\bigO(N)$ kriging evaluation. Compare the original
full-$N$ joint GN warm-up, which cost $22$\,s at $N=800$ and $114$\,s at
$N=1600$ (an $\bigO(N^{\sim 2})$ trajectory that would reach $\sim\!1800$\,s at
$N=6400$). (ii) The total grows essentially linearly (quantified below).
(iii) $\mathrm{rel}_a$ is flat-to-improving
in $N$, confirming robust two-mode recovery at every size; the warm-up's coarse
solve drives the coarse loss far below the full-$N$ warm-up's, so the kriged
seed is actually a better basin point (the coarse-warm-up $\mathrm{rel}_a$
\emph{improves} from $7.89\%$ to $7.40\%$ as $N$ grows, beating the full-$N$
warm-up's $\approx 13\%$). Concretely, a $16\times$ increase in $N$ ($800\to
12800$) raises the total wall by only $19.8\times$ --- an overall empirical
exponent $\approx 1.08$ --- and the per-doubling ratio settles toward $\approx
2.1$ at the largest sizes (the asymptotic regime where the near-linear terms
dominate the fixed coarse-warm-up overhead).

\paragraph{Per-step $u$-iteration counts.} Across the whole sweep the forward
pCG inside \eqref{eq:bordered} converges in a number of iterations that does not
grow with $N$ (high-$60$s to low-$70$s), and the $w$-update in $\sim\!13$
iterations --- the hallmark of the screening preconditioner acting on a
well-conditioned covariance operator, exactly as in the forward solver.

\subsection{Isolating the field scaling: fixed $N_d$}\label{sec:fixedNd}

Table~\ref{tab:scaling} grows the data set with the field, $N_d\propto N$, so it
folds the dense $\bigO(N_d^2)$--$\bigO(N_d^3)$ data-block term into the measured
exponent. To isolate the pure field/PDE scaling --- the part that should be
near-linear --- we repeat the sweep with the number of noisy measurements
\emph{held fixed} at $N_d=300$ while $N$ grows $16\times$ (Table~\ref{tab:fixedNd}).

\begin{table}[h]
\centering
\caption{Fixed-$N_d$ sweep ($N_d=300$ for all $N$). With the data block frozen,
the remaining cost is the field work (factor builds, forward pCG, kriging) and is
cleanly near-linear; coefficient recovery is unchanged.}
\label{tab:fixedNd}
\begin{tabular}{rrrrrr}
\toprule
$N$ & $N_d$ & Stage-1 (s) & Stage-2 (s) & Total (s) & $\mathrm{rel}_a$ \\
\midrule
$800$   & $300$ & $6.9$  & $9.8$   & $23.4$  & $7.89\%$ \\
$1600$  & $300$ & $7.0$  & $21.2$  & $39.8$  & $7.81\%$ \\
$3200$  & $300$ & $7.6$  & $67.5$  & $96.8$  & $7.73\%$ \\
$6400$  & $300$ & $11.1$ & $149.8$ & $201.1$ & $7.62\%$ \\
$12800$ & $300$ & $15.8$ & $331.3$ & $423.5$ & $7.40\%$ \\
\bottomrule
\end{tabular}
\end{table}

\noindent A $16\times$ increase in $N$ raises the total by $18.1\times$ (overall
empirical exponent $\approx 1.05$; top-end per-doubling ratio $2.11$). The
revealing comparison is against Table~\ref{tab:scaling} at $N=12800$: there
$N_d$ had grown to $4800$ ($16\times$, so the dense block's $N_d^3$ grew
$\sim\!4000\times$), yet the total wall is only $\sim\!6\%$ larger ($448.9$ vs.\
$423.5$\,s) and the exponent barely moves ($1.08$ vs.\ $1.05$). \emph{The dense
data block is not the bottleneck in this range:} the cost is dominated by the
per-step field operations --- the screened factor builds and forward triangular
solves --- which scale near-linearly. Holding $N_d$ fixed does not visibly
``clean up'' the curve precisely because the peeled-off data block was already
sub-dominant; the trailing-block identity \eqref{eq:trailing} kept the only
field-coupled data cost at $\bigO(N)$, leaving the small dense $\bigO(N_d^3)$
factorization as a cheap, $N$-decoupled additive constant.

\subsection{Verification of the precision-factorization claims}\label{sec:verif}

Two precision-factorization claims underlie the solver design: that incomplete
Cholesky \emph{does} factor a screened precision-plus-nugget $\Th\inv+R\inv$
(justifying the noisy-regression machinery), but \emph{does not} factor the
Jacobian-sandwiched GN Hessian $H$ (justifying why the warm-up sidesteps $H$,
Sec.~\ref{sec:warmup}). We checked both directly; the conclusions are distilled
as Insight~6 (Sec.~\ref{sec:insights}).

\paragraph{(a) ichol of $\Th\inv+R\inv$ works (at moderate $\sigma$).} We built
the noiseless KL factor of a Matérn-$7/2$, $\ell=1$, $\rho=3$ covariance $\Th$ on
$N$ random points, ran Alg.~4.1 to obtain $\tilde U\tilde U\T\approx\Th\inv+
\sigma^{-2}I$, and preconditioned CG (rtol $10^{-8}$) on $\Sigma=\Th+\sigma^2 I$.

\begin{table}[h]
\centering
\caption{ichol (Alg.~4.1) as a preconditioner for $\Sigma=\Th+\sigma^2 I$. At
moderate $\sigma^2$ the iteration count is essentially $N$-independent and the
factor stays sparse; at the tiny $\sigma$ of the Darcy problem it stalls.}
\label{tab:ichol}
\begin{tabular}{rrrr}
\toprule
& \multicolumn{1}{c}{$\sigma^2=10^{-2}$} & & \multicolumn{1}{c}{$\sigma^2=10^{-6}$}\\
\cmidrule(lr){2-2}\cmidrule(lr){4-4}
$N$ & CG iters ($\mathrm{nnz}(\tilde U)/\mathrm{nnz}(U)$) & & CG iters \\
\midrule
$400$  & $20$ \,($1.5$)  & & no conv.\ ($500^{+}$) \\
$800$  & $35$ \,($1.7$)  & & no conv.\ ($500^{+}$) \\
$1600$ & $34$ \,($1.8$)  & & no conv.\ ($500^{+}$) \\
$3200$ & $42$ \,($1.8$)  & & no conv.\ ($500^{+}$) \\
$6400$ & $46$ \,($1.9$)  & & no conv.\ ($500^{+}$) \\
\bottomrule
\end{tabular}
\end{table}

\noindent At moderate noise this is a genuine, near-$N$-independent precision
preconditioner. As $\sigma\to0$ (our $\sigma=10^{-3}$, i.e.\ $\sigma^2=10^{-6}$)
the nugget $R\inv$ dominates and the factored correction degenerates: CG does not
converge in $500$ iterations at any $N$. That stalling, plus the structural fact
that Alg.~4.1 requires a \emph{uniform positive} nugget on every diagonal (it
cannot encode the exact, $R=0$, PDE/boundary rows of the $u$-step), is exactly
why the solver peels the data block (Sec.~\ref{sec:alt}) instead of factoring
through the noise.

\paragraph{(b) incomplete factorization of $H$ fails.} We dumped the fine joint
Hessian $H$ at the first GN step and preconditioned CG on $Hx=g$ with an
incomplete LU (the closest available drop-tolerance factorization) at three fill
levels, alongside plain and Jacobi-preconditioned CG (rtol $10^{-6}$,
maxiter $3000$).

\begin{table}[h]
\centering
\caption{Preconditioned CG iteration count on the joint GN Hessian $Hx=g$. No
incomplete factorization converges at any $N$ or fill; only the \emph{exact}
block factorization (used in the warm-up, at $\bigO(N^{2.8})$ cost) solves it.}
\label{tab:Hfail}
\begin{tabular}{rrrrc}
\toprule
$N$ & $\kappa(H)$ & plain CG & Jacobi CG & spilu \\
\midrule
$400$  & $2.0\times10^{9}$ & $3000^{+}$ & $2653$ & no conv. \\
$800$  & $5.2\times10^{8}$ & $3000^{+}$ & $2409$ & no conv. \\
$1600$ & $3.7\times10^{8}$ & $3000^{+}$ & $2122$ & no conv. \\
\bottomrule
\end{tabular}\\[2pt]
{\footnotesize spilu tested at three drop/fill settings up to
$\mathrm{nnz}(\text{factor})=2.5\times\mathrm{nnz}(H)$; none converged in $3000$
iterations. ``$3000^{+}$'' = hit the iteration cap.}
\end{table}

\noindent Even at $\mathrm{nnz}(\text{factor})=2.5\times\mathrm{nnz}(H)$ the
incomplete factorization gives no convergence --- the sandwiched-operator
obstruction, distinct from the well-behaved precision-plus-nugget of~(a). This is
the quantitative basis for sidestepping $H$ via the coarse warm-up.

% =========================================================================
\section{Key insights: inverse problems with several GP priors and noise
nuggets}\label{sec:insights}
% =========================================================================

The Darcy problem is a concrete instance of a recurring structure: an inverse
problem with \emph{more than one} Gaussian-process prior (here $u$ and $w$) and
one or more \emph{noise nuggets} --- $\sigma^2 I$ regularizations attached to
noisy or soft observation functionals. The following lessons are what make the
sparse-Cholesky machinery deliver forward-solve efficiency in that structure;
they are stated so as to transfer to other multi-field GP-collocation inverse
problems.

\paragraph{Insight 1 --- one screened factor per field; never sum kernels.}
With $k$ independent priors $\phi_m\sim\cGP(0,K_m)$ the loss is a sum of $k$
RKHS norms $\sum_m \xi_m\T\Th_m\inv\xi_m$ \eqref{eq:loss}, \emph{each a
single-kernel covariance}. Keep them that way: build one KL factor per kernel and
apply them blockwise. The temptation is to collapse the fields into one composite
observation functional $L=\sum_m c_m L_m$ acting on a summed kernel
$\sum_m c_m^2 K_m$; this destroys screening, because the slow features of the
smoothest kernel (here the long-range $\Delta u$ mode of $\ell_u=1$) get mixed
into every entry and the truncated factor needs a large radius $\rho$ to stay
accurate. Separateness is what keeps each $\rho$ in the elliptic regime
($\rho=3$).

\paragraph{Insight 2 --- give each field its own length scale (screening budget).}
Different unknowns in an inverse problem have different regularity, and the
sparsity radius needed for accurate screening is set \emph{per field}. The state
$u$ is smooth (it solves an elliptic PDE) and is given a long $\ell_u$ so its
derivatives, which pin the coefficient through the PDE, stay clean; the
coefficient $w$ is rough and is given a short $\ell_w$ so its factor stays local.
A single shared length scale forces a compromise that flattens the recovered
coefficient. The multi-GP formulation is what \emph{allows} this per-field tuning
--- another reason not to merge the priors.

\paragraph{Insight 3 --- a noise nugget is benign in isolation but stiff when
co-located with a hard constraint.} Adding $\sigma^2 I$ to a covariance block is
a regularization and normally \emph{improves} conditioning. The danger in the
inverse setting is geometric: when a \emph{noisy} functional sits on (or very
near) a point already carrying a \emph{noiseless/hard} functional --- here the
data $\delta_{x_j^d}$ co-located with the PDE collocation point $x_j^\Omega$, with
tiny $\sigma$ --- the noisy row is nearly redundant with the hard row. Its
conditional variance given the hard block is $\approx\sigma^2$, so the augmented
Gram acquires a tight cluster of near-null modes whose count grows with the
number of nuggets $N_d$. A local screened factor of the \emph{augmented} operator
cannot resolve them at any practical $\rho$ (raising $\rho$ fixes one problem
size but the stiff-mode count scales $\propto N_d$). \emph{Diagnose this by the
ratio of nugget magnitude to the conditional variance of the noisy functional
under the hard constraints; when $\sigma$ is small relative to that variance, do
not precondition through the nugget block.}

\paragraph{Insight 4 --- peel the nugget block with an exact Schur step; never
fold it into the preconditioner.} Partition the functionals into the hard
forward block $F$ (which screens perfectly and reproduces the forward solve) and
the soft/nuggeted block $d$. Solve by the exact bordered/Schur formulas
\eqref{eq:bordered}: apply $\Th_{FF}\inv$ with the forward pCG ($N$-independent
iterations, the field's own small factor as preconditioner), couple through
single big-factor mat-vecs $\Th_{Fd}$, and reduce to a small dense
$N_d\times N_d$ system $S=S_0+\sigma^2 I$. This keeps the well-conditioned
forward block on its native conditioning and confines the stiffness to the small
dense block where it is harmless.

\paragraph{Insight 5 --- read the nugget-block posterior covariance off the
trailing precision block, not from forward solves.} The one object the Schur
step needs, the data-block posterior precision $S_0\inv=(\Thtr\inv)_{dd}$, is
\emph{already present} in the joint KL factor: by \eqref{eq:trailing} it is the
dense Gram $U_d U_d\T$ of the $N_d$ permuted nugget rows of $U$, formed in
$\bigO(\mathrm{nnz}(U_d))=\bigO(N)$ time with no forward solves. This is the
identity that turns an $\bigO(N\cdot N_d)$ coupling into a near-linear one and is
the linchpin of the whole scheme. (A single spurious large-variance mode appears
in the screened $S_0$; it is benign, merely returning that data direction to the
accurate forward solve.) The fixed-$N_d$ experiment (Sec.~\ref{sec:fixedNd})
confirms the payoff: with the field-coupled data cost reduced to $\bigO(N)$, the
remaining dense $\bigO(N_d^3)$ factorization is an $N$-decoupled constant and the
total scales near-linearly.

\paragraph{Insight 6 --- factor a precision only when it is a screened
covariance-precision plus a diagonal; never the Jacobian-sandwiched normal form.}
``Sparse Cholesky belongs on the covariance, not the precision'' is too blunt:
incomplete Cholesky factors a \emph{precision} just fine when that precision is
$\Th\inv+R\inv$ --- a screened covariance-precision plus a diagonal nugget. This
is precisely the Schäfer--Katzfuss--Owhadi noisy variant (Alg.~4.1): build the
noiseless KL factor $U$ of the covariance $\Th$ ($UU\T\approx\Th\inv$), then a
\emph{second} incomplete Cholesky on the $UU\T$ pattern gives
$\tilde U\tilde U\T\approx\Th\inv+R\inv$. As a preconditioner it solves
$\Sigma=\Th+\sigma^2 I$ in an $N$-robust $\sim\!20$--$46$ CG iterations at
moderate $\sigma^2$ (verified, Sec.~\ref{sec:verif}). The object that genuinely
resists factorization is the GN normal-equations Hessian
$H=J_v\T\Th_u\inv J_v+J_w\T\Th_w\inv J_w+D$ \eqref{eq:H}: the Jacobian
sandwiching makes it a squared, biharmonic-like operator whose exact factor fills
in ($\bigO(N^{2.8})$) and whose ill-conditioning ($\kappa(H)\sim10^{9}$) is
\emph{not} of the screened form; an incomplete factorization of $H$ does not
converge at any $N$ or fill, and Alg.~4.1 does not even apply since
$H\ne\Th\inv+\mathrm{diag}$. Two further caveats sharpen the rule, both
verified in Sec.~\ref{sec:verif}: (i) even the legitimate ichol of
$\Th\inv+R\inv$ \emph{degrades as $\sigma\to0$} --- at our $\sigma=10^{-3}$ it
stalls --- and (ii) Alg.~4.1 needs a \emph{uniform positive} nugget on every
diagonal, so it cannot represent the exact ($R=0$) PDE/boundary rows of the
$u$-step. Hence we peel the data block (Insight~4) rather than ichol through the
noise, and sidestep $H$ entirely (Insight~7).

\paragraph{Insight 7 --- decouple the nonconvexity: only the basin-defining field
needs the expensive nonlinear solve.} Multi-field inverse problems are non-convex
through the field coupling \eqref{eq:v3}. But the expensive global step (joint GN
on $H$) is needed \emph{only} to land in the right basin, and the alternating
refinement re-solves each field from scratch given the other. So run the
nonlinear warm-up on a \emph{fixed coarse subset} ($\bigO(1)$ in $N$) and krige
its output field onto all $N$ points ($\bigO(N)$) --- the basin is a property of
the coefficient, captured at coarse resolution. This removes the only inherently
global, superlinear computation from the fine-$N$ path. (It is not multigrid: one
coarse solve, one kriging evaluation, no hierarchy or cycles.)

% =========================================================================
\section{Discussion}\label{sec:discussion}
% =========================================================================

The solver attains forward-elliptic efficiency for a two-field nonlinear inverse
problem by keeping \emph{every} linear solve in covariance form and never
touching the precision Hessian. The single remaining super-linear term is the
dense $N_d\times N_d$ nugget-block solve; the fixed-$N_d$ sweep
(Sec.~\ref{sec:fixedNd}) shows it is sub-dominant in the tested range, and
screening $S$ (whose posterior covariance decays in space because the nugget
sites are spatially distributed) would close the last gap and is the natural next
step. As implemented, the method already delivers near-linear total cost with
$N$-robust two-mode recovery over $N=800$--$12800$.

\paragraph{Reproduction.} \texttt{examples/darcy\_inverse\_hybrid.py} with
\texttt{-{}-warmup-mode coarse} (default), \texttt{-{}-noise-mode data-bordered},
\texttt{-{}-kernel-sigma-u 1.0 -{}-kernel-sigma-w 0.2 -{}-rho-small 3
-{}-warmup-steps 6 -{}-alt-steps 3}, sweeping \texttt{-{}-N-domain} /
\texttt{-{}-N-data}.

\end{document}
